\section{一、引言}\label{ux4e00ux5f15ux8a00}

数据要素在企业价值创造中的作用持续上升，但会计计量体系仍然滞后（洪永淼和史九领，2024）。2022年``数据二十条''和2024年《企业数据资源相关会计处理暂行规定》在制度层面迈出了重要一步，然而现实困难依然突出：历史成本模式难以反映数据资源的动态价值特征，商业模式差异又使同一数据资源在不同场景下价值迥异（黄世忠等，2023；Mohan等，2026）。这一计量缺口意味着投资者难以仅凭财务报表完整评估企业价值，资本市场对相关企业的定价效率因此受到影响。

企业在年报中披露的数据要素利用信息可能通过两条相反的路径影响定价效率。从信息供给的角度看，数据利用的实质性描述降低了投资者面临的信息不对称，有助于加速价格发现（Diamond和Verrecchia，1991）。但从信息模糊的角度看，数据驱动商业模式的内在复杂性和估值标准的缺失可能加剧不确定性，反而阻碍价格发现（Zhang，2006）。净效果取决于哪种力量占主导，需要经验数据加以识别。中国A股市场的年报经过审计、受信息披露监管约束，数据要素政策框架也在逐步完善，这些条件可能有利于信息供给效应的发挥，但仍需实证检验。

既有研究主要沿两条脉络展开。第一条脉络关注数据要素对企业行为的影响。朱康和唐勇（2025）发现数据要素利用降低了金融资产配置并促进实体投资，李姝等（2025）提供了数据资产缓解融资约束的互补证据，Shao等（2024）从经济政策不确定性视角报告了类似发现。这些研究揭示了数据利用在企业决策层面的实质效应，但研究视角大多停留在企业行为层面，尚未系统检验资本市场如何吸收这类信息。

第二条脉络直接考察信息环境与定价效率的关系。会计信息质量的改善能够降低股价延迟（Callen等，2013），非财务信息的前瞻性有助于提升定价效率（Zhang和Wang，2024），另类数据对财务预测的改善在较长时间跨度上更为明显（Dessaint等，2024）。在中国市场，行业信息披露指引（Gao等，2024）和主流财经媒体关注（Yin和Hu，2024）均被发现能改善股价信息含量。与数据要素更直接相关的研究中，数字化转型降低股价延迟的初步证据已被报告（Feng等，2024），数据资产信息披露降低股价同步性和股价延迟的结论也在特定样本中得到支持（Sun和Du，2024；Chen，2024；Huang等，2026；Wei等，2025；李世刚等，2025）。

然而，``数字化转型''\,``数据资产披露''与``数据要素利用''是相关但不等同的概念。数字化转型侧重技术采纳，数据资产披露侧重会计确认，数据要素利用关注的是企业对数据资源的实际应用。本文后续的构念边界检验也表明，这三者并非可以被完全机械切分：数据利用测度在剔除泛化词和控制篇幅后仍保留解释力，但与更宽泛数字叙事之间仍存在交叠。因此，本文的目标不是将其表述为与数字化转型完全正交的概念，而是更审慎地识别其相对更贴近``数据资源实际利用''的部分。

如果年报中的数据利用信息确实帮助市场更快理解企业价值，那么这种作用不应只体现在总量上的股价延迟下降，还应体现在更细的传导环节上。其一，信息披露可能降低分析师对企业基本面的认知分歧，使预测更趋一致。其二，数据驱动的经营优化可能降低企业现金流波动，从而减少估值不确定性。其三，数据能力可能帮助企业优化客户和供应商匹配，缓解供应链集中度过高带来的基本面脆弱性。这三条路径分别对应信息质量、经营稳定性和供应链韧性，能够帮助识别数据要素利用如何从年报文本传导至股票价格。

本文旨在构建这一完整路径。与朱康和唐勇（2025）聚焦企业决策行为和李世刚等（2025）以入表口径度量数据资产化不同，本文直接考察年报中数据要素实际应用的实质性描述如何影响股价的信息吸收速度，并系统检验其与分析师预测分歧、现金流波动和供应链集中度之间的关系。

本文以2011至2024年沪深A股上市公司为样本，基于年报全文考察数据要素利用信息对股价延迟的影响。度量上采用两种互补策略：关键词频率指标扫描年报中与数据收集、分析、应用和价值变现相关的表述，捕捉利用信息的广度；大语言模型对关键词上下文进行语义评分，刻画利用信息的深度。本文将DU\_kw作为基准处理变量，将DU\_llm作为语义增强测度，并进一步构造长度标准化的DU\_llm\_lenstd用于检验语义深度是否在控制篇幅后仍保留解释力。与此同时，本文进一步利用现有五维词典和年报章节特征，构造数据利用链条完整度（DUchain\_count）、闭环式表述（DUclosedloop）、窄口径核心强度（DUcore）和MD\&A章节限定口径（DUkw\_mda），以检验测度能否从泛化数字叙事进一步推进到更接近数据资源实际使用的口径。研究发现，年报中的数据要素利用信息显著降低股价延迟；进一步以前一期DU解释当期PriceDelay，并结合工具变量、双重机器学习、Oster界检验、Heckman两阶段法和安慰剂检验，结论保持稳定，多种识别练习支持因果解释。构念边界检验显示，核心结果在篇幅控制和去泛化词口径下仍然成立，但GenericNarr联合回归表明二者存在较强重叠，因此概念独立性应作审慎表述。增强测度结果显示，更完整的数据利用链条、闭环式表述以及MD\&A中的相关披露同样与更低的股价延迟相关，但这些口径在与DU\_kw或GenericNarr联合进入时通常不能稳定替代基准处理变量，因此更适合作为支持性测度。机制结果显示，数据要素利用信息与更低的现金流波动和更低的供应链集中度稳健相关，并在分析师预测分歧维度上呈现补充性证据。扩展分析进一步表明，关键词广度与语义深度的错配会加剧股价延迟，而语义深度在与关键词广度同时进入时仍保留额外解释力。异质性结果进一步表明，该效应不仅在2022年之前更强，也在高分析师关注企业和非主板企业中表现得更为明显；此外，行业整体披露水平对企业股价延迟存在独立的负向作用，绝对量级在滞后规格下甚至超过企业自身披露，体现出行业层面的信息环境同群溢出。

本文的贡献主要体现在以下三个方面。第一，构建了关键词频率与大语言模型语义评分的双测度体系，并进一步通过去泛化词、联合回归、五维闭环度、闭环式表述和MD\&A章节限定口径检验测度边界。本文据此不再将``数据要素利用''表述为与更宽泛数字叙事完全分离的概念，而是更审慎地识别其相对更贴近实际数据利用的部分；其中，闭环度和经营章节口径为``更完整、更嵌入''的数据利用信息提供了支持性证据。第二，本文将数据要素研究从企业行为后果推进到资本市场的信息效率后果，并从信息分歧、经营稳定性与供应链韧性三个维度识别传导路径；进一步的WashGap分析显示，市场更倾向于奖励具有实质内容和经营嵌入的数据利用披露，而非单纯广泛但浅层的概念性表述。第三，本文将时间顺序更清晰的滞后规格前置，并结合工具变量、双重机器学习、Oster界检验、Heckman选择修正和安慰剂检验，形成多层次的识别框架，以更审慎的方式支持因果解释。

本文余下部分安排如下。第二部分建立理论框架并导出研究假说。第三部分介绍研究设计。第四部分依次报告基准回归、构念边界与测度补强、时间顺序与内生性识别、双重机器学习估计和稳健性检验。第五部分进一步检验传导机制与广度---深度错配。第六部分报告政策阶段、信息中介活跃度和上市板块异质性分析。第七部分总结全文并讨论政策启示。
